% /solutions/conference-talks/conference-ornate-20min.fr.tex, 22/02/2006 De Sousa
\documentclass{beamer}
\usepackage{listings}
\usepackage{mdframed}
\usepackage{tikz}
% Ce fichier est un exemple d'expos\'e

% - pour des conf\'erences,
% - d'une dur\'ee approximative de 20 minutes,
% - avec un style ornemental.


% Copyright 2004 by Till Tantau <tantau@users.sourceforge.net>.
%
% Traduction de Philippe De Sousa <philippejjg@free.fr>
%
% En principe, ce fichier peut être redistribu\'e et/ou modifi\'e conform\'ement
% aux termes de la GNU Public License, version 2.
%
% Cependant, ce fichier est suppos\'e comme \'etant un "exemple-type" qui peut être modifi\'e
% selon vos propres besoins. Pour cette raison, si vous utilisez ce fichier en tant qu'
% "exemple-type" et non sp\'ecifiquement pour le distribuer en tant que partie d'un
% package ou programme, je vous donne la permission exceptionnelle de copier librement et
% de modifier ce fichier et même d'effacer ce message de copyright.
\usepackage{pdfpages}



\usepackage[francais]{babel}
% or autre comme par exemple \usepackage[english]{babel}

\usepackage[utf8]{inputenc}
% or autre

\usepackage{float}
\usepackage{graphicx}
\usepackage{wrapfig}
\usepackage{times}
\usepackage[T1]{fontenc}
% Or autre. Notez que le codage et la fonte doivent être assortis. Si T1
% ne paraît pas tr\`es esth\'etique, essayer d'effacer la ligne contenant fontenc.
\makeatletter
\hypersetup{pdfpagemode=FullScreen}

\mode<presentation> {
  \usetheme{UNLTheme}
  % ou autre ...
	
 	% \setbeamercovered{none}
  % ou autre chose (il est \'egalement possible de supprimer cette ligne)
}

\title[] % (facultatif, \`a utiliser uniquement si le titre de l'article est trop long)
{Logiciel de devis et Facturation}
\subtitle {FactDev}

\author[
\textbf{F}lorent\\
\textbf{A}ntoine\\
\textbf{C}édric\\
Manan\textbf{T}soa
] % (facultatif, \`a utiliser seulement avec plusieurs auteurs)
{Antoine de \bsc{Roquemaurel}\newline Florent \bsc{Berbie}\newline Cédric \bsc{Rohaut}\newline Manantsoa Andriamihary \bsc{Razanajatovo}}

% - Composez les noms dans l'ordre dans lequel ils apparaîtrons dans l'article
% - Utilisez la commande \inst{?} uniquement si les auteurs ont des affiliations
%   diff\'erentes.

\institute[] % (facultatif mais g\'en\'eralement n\'ecessaire)
{
  Universit\'e Toulouse III -- Paul Sabatier \\
  M1 Informatique -- Développement Logiciel 
  \vspace{-10px}
}
% - Utilisez la commande \inst uniquement s'il y a plusieurs affectations
% - Faîtes quelque chose de simple, personne ne s'int\'eresse \`a votre adresse

\date[ ~ ~ ~ JJ / MM / 2015] % (facultatif, peut être une abr\'eviation du nom de la conf\'erence)
{JJ / MM / 2015}
% - Utilisez \`a la fois le nom de la conf\'erence et son abr\'eviation.
% - N'a pas r\'eellement d'importance pour l'assistance qui sera pr\'esente lors de la conf\'erence,
%   mais en a surtout pour les personnes (y compris vous-même) qui liront les transparents en ligne.

\subject{Plateforme de tests automatis\'es}
% Ins\'er\'e uniquement dans la page d'information du fichier PDF. Peut être
% supprim\'e.


% Si vous avez un fichier nomm\'e "universit\'e-logo-nomfichier.xxx", où xxx
% est un format graphique accept\'e par latex ou pdflatex (comme par exemple .png),
% alors vous pouvez ins\'erer votre logo ainsi :

 \pgfdeclareimage[width=2.5cm]{le-logo}{../../images/FACT_official.png}
 \logo{\pgfuseimage{le-logo}}


% \`a supprimer si vous ne voulez pas que la table des mati\`eres apparaisse
% au d\'ebut de chaque sous-section : 
\AtBeginSection[] {
  \begin{frame}<beamer>{Ligne directrice}
    \tableofcontents[currentsection]
  \end{frame}
}
% Si vous souhaitez recouvrir vos transparents un \`a un,
% utilisez la commande suivante (pour plus d'info, voir la page 74 du manuel
% d'utilisation de Beamer (version 3.06) par Till Tantau) :

%\beamerdefaultoverlayspecification{<+->}


\setbeamertemplate{footline}{
	\leavevmode%
	\hbox{\hspace*{-0.6cm}
	\begin{beamercolorbox}[wd=.19\paperwidth,ht=2.25ex,dp=1ex,center]{section in head/foot}%
		\usebeamerfont{section in head/foot}\insertshortdate{}
	\end{beamercolorbox}%

	\begin{beamercolorbox}[wd=.76\paperwidth,ht=2.25ex,dp=1ex,center]{section in head/foot}%
		\insertprogressbar
		\usebeamerfont{section in head/foot} Développement d'un logiciel de facturation : FactDev 
		\end{beamercolorbox}%

		\begin{beamercolorbox}[wd=.10\paperwidth,ht=2.27ex,dp=1ex,right]{section in head/foot}%
			\usebeamerfont{section in head/foot}\hspace*{2em}
			\hspace{-10px}\textbf{\insertframenumber{}} / \inserttotalframenumber\hspace*{2ex}
		\end{beamercolorbox}}%
		\vskip0pt%
		}

		\setbeamertemplate{navigation symbols}{}
		\begin{document}
		\begin{frame}
			\titlepage
		\end{frame}
		\begin{frame}{Les parties prenantes}
			\begin{itemize}
				\item \'Equipe FACT
					\begin{itemize}
						\item {Antoine de \bsc{Roquemaurel}
						\item Florent \bsc{Berbie}
						\item Cédric \bsc{Rohaut}
						\item Manantsoa Andriamihary \bsc{Razanajatovo}}
					\end{itemize}
					\vfill
				\item Encadrant
					\begin{itemize}
						\item Frédéric \bsc{Migeon}
					\end{itemize}
					\vfill
				\item Client
					\begin{itemize}
						\item Antoine de \bsc{Roquemaurel}
					\end{itemize}
					\vfill
			\end{itemize}
		\end{frame}
		\begin{frame}
			\tableofcontents	
		\end{frame}
\section{Le projet}
\subsection{FactDev ?}
\begin{frame}{FactDev : Un logiciel de Devis et Facturation}
\begin{itemize}
	\item Usage professionnel
	\item Gestion de 
		\begin{itemize}
			\item Clients
			\item Projets associés
			\item Tarifs
			\item Génération de documents 
			\item Recherche
		\end{itemize}
	\item Développement itératif
\end{itemize}
\end{frame}
\subsection{Les outils}
\begin{frame}{Les outils utilisés}
	\begin{figure}[H]
		\centering
	\includegraphics[height=1.5cm]{logos/qt.png}~~
	\Huge \LaTeX

	\includegraphics[height=1.5cm]{logos/travis.png}~~
	\includegraphics[height=1.3cm]{logos/coveralls.png}

	\includegraphics[height=1.5cm]{logos/git.png}~~
	\includegraphics[height=1.2cm]{logos/github.png}

	\includegraphics[height=0.8cm]{logos/irc.png}~~
	\includegraphics[height=0.5cm]{logos/doxygen.png}~~ 
	\includegraphics[height=0.8cm]{logos/sqlite.png}
	\caption{Les différents outils}
\end{figure}
	%% Valgrind ? 
	
\end{frame}
\section{Méthodologie}
\subsection{Scrum}
\begin{frame}{Scrum, une méthodologie de développement}
	\begin{itemize}
		\item 
	\end{itemize}<++>
	%% TODO add figure Sprints/user stories
\end{frame}
\subsection{Revue de code}
\begin{frame}
%% Pull Requests	
\end{frame}
\subsection{Intégration continue}
\begin{frame}
% Travis	
% Couverture, build, test, PR
\end{frame}
\section{Les résultats}
\subsection{Pour la méthode ?}
\begin{frame}
% xx issues, xx PR, xx commits, xx branches, xx builds, xx de couverture, …
\end{frame}

\subsection{Et le logiciel ?}
\begin{frame}
% Screens 
\end{frame}

\section*{Conclusion}
\begin{frame}
% It was cool. But difficult. But we learn. And it was cool. And ? 
	
\end{frame}

% Slide for questions
\begin{frame}{Avez-vous des questions ?}
	\begin{figure}[H]
		\centering
		\includegraphics[width=5cm]{interrogation.png}
	\end{figure}
\end{frame}
\end{document}
